\chapter{Automatic Differentiation}
\label{chap:autodiff}

This chapter specifies how \lupnt{} obtains derivatives. Every
state-transition matrix in \cref{chap:dynamics} and \cref{chap:integration},
every measurement Jacobian consumed by the estimators of \cref{chap:filters},
and every force-model sensitivity used in parameter estimation is produced by
\emph{forward-mode automatic differentiation} (AD) of the model code itself
--- not by hand-coded partial derivatives and not by finite differences. The
consequence is a hard contract on how models may be written: a model that
respects the rules of \cref{sec:ad-rules} gets exact first derivatives for
free, and a model that violates them silently returns wrong ones.

The differentiable scalar and the imported AD entry points are declared in
\file{cpp/lupnt/core/definitions.h}; the parallel Jacobian driver is in
\file{cpp/lupnt/numerics/math_utils.cc}; the filter-facing wrappers are in
\file{cpp/lupnt/numerics/filters/filter_utils.cc}. The underlying library is
\code{autodiff} \citep{leal2018autodiff}, whose containers are Eigen matrices
\citep{eigen2010}. The theory --- forward and reverse accumulation, seeding,
and the cost model --- follows the standard treatment of algorithmic
differentiation \citep{griewank2008}.

\section{The Differentiable Scalar \texorpdfstring{\cpp{Real}}{Real}}
\label{sec:ad-real}

\subsection{Definition}

\lupnt{}'s fundamental arithmetic type is
\begin{equation}
  \label{eq:ad-real}
  \cpp{Real} \equiv \code{autodiff::real} \equiv \code{autodiff::Real<1, double>},
\end{equation}
a first-order truncated Taylor number over \code{double}. Concretely it is a
two-element array: slot $0$ holds the value, slot $1$ holds one directional
derivative. Writing a \cpp{Real} as the pair
\begin{equation}
  \label{eq:ad-dual}
  \tilde{x} = \big(x,\ \dot{x}\big),
  \qquad x = \tilde{x}.\code{val()},
  \qquad \dot{x} = \tilde{x}[1],
\end{equation}
the type is algebraically the dual numbers $\mathbb{R}[\epsilon]/(\epsilon^2)$
under the identification $\tilde{x} = x + \dot{x}\,\epsilon$ with
$\epsilon^2 = 0$.

\begin{lstlisting}[language=cpp, caption={Storage layout of the differentiable scalar. Source: \code{autodiff/forward/real/real.hpp} (vendored dependency)}, label={lst:ad-realtype}]
// autodiff/forward/real/real.hpp :: Real
template <size_t N, typename T>
class Real {
  std::array<T, N + 1> m_data = {};
public:
  constexpr auto val() -> T& { return m_data[0]; }
  constexpr auto val() const -> const T& { return m_data[0]; }
  constexpr auto operator[](size_t i) -> T& { return m_data[i]; }
  constexpr auto operator[](size_t i) const -> const T& { return m_data[i]; }
};

using real1st = Real<1, double>;
using real    = real1st;    // == lupnt::Real
\end{lstlisting}

\subsection{The arithmetic}

Every elementary operation is overloaded to propagate the second slot by the
chain rule. For $\tilde{x} = (x,\dot x)$ and $\tilde{y} = (y,\dot y)$,
\begin{align}
  \label{eq:ad-arith}
  \tilde{x} \pm \tilde{y} &= \big(x \pm y,\ \dot x \pm \dot y\big), &
  \tilde{x}\,\tilde{y}    &= \big(xy,\ x\dot y + \dot x y\big), \nonumber\\
  \frac{\tilde{x}}{\tilde{y}}
    &= \left(\frac{x}{y},\ \frac{\dot x y - x \dot y}{y^{2}}\right), &
  \sqrt{\tilde{x}} &= \left(\sqrt{x},\ \frac{\dot x}{2\sqrt{x}}\right),
\end{align}
and for a general differentiable elementary function $g$,
\begin{equation}
  \label{eq:ad-chain}
  g(\tilde{x}) = \big(g(x),\ g'(x)\,\dot{x}\big).
\end{equation}
\Cref{eq:ad-chain} is the whole mechanism. Two of its instances matter
repeatedly below because their derivative factor is unbounded:
\begin{equation}
  \label{eq:ad-acos}
  \arccos(\tilde{x})
    = \left(\arccos x,\ \frac{-\dot{x}}{\sqrt{1-x^{2}}}\right),
  \qquad
  \arcsin(\tilde{x})
    = \left(\arcsin x,\ \frac{\dot{x}}{\sqrt{1-x^{2}}}\right),
\end{equation}
both of which diverge as $|x| \to 1$ and evaluate to $0/0 = \mathrm{NaN}$ at
$|x| = 1$ exactly. \Cref{sec:ad-rules} specifies how models must handle this.

\subsection{Containers}

Every \lupnt{} vector and matrix type is an Eigen container templated on
\cpp{Real} \citep{eigen2010}, with a plain-\cpp{double} counterpart --- suffix
\code{d} --- used for storage, input/output, and the interfaces that cross
into non-differentiated code such as the filters' covariance algebra.

\begin{lstlisting}[language=cpp, caption={Differentiable scalar, AD entry points, and container aliases. Source: \code{cpp/lupnt/core/definitions.h}}, label={lst:ad-defs}]
// cpp/lupnt/core/definitions.h
namespace ad = autodiff;

using Real = ad::real;      // forward mode, first order
using ad::at;
using ad::derivative;
using ad::jacobian;
using ad::wrt;

template <int rows, int cols> using Mat  = Matrix<Real,   rows, cols>;
template <int rows, int cols> using Matd = Matrix<double, rows, cols>;
template <int size>           using Vec  = Matrix<Real,   size, 1>;
template <int size>           using Vecd = Matrix<double, size, 1>;

// ... expanded by DEFINE_VECTORS_MATRICES():
using VecX  = Matrix<Real,   Eigen::Dynamic, 1>;
using VecXd = Matrix<double, Eigen::Dynamic, 1>;
using MatX  = Matrix<Real,   Eigen::Dynamic, Eigen::Dynamic>;
using MatXd = Matrix<double, Eigen::Dynamic, Eigen::Dynamic>;
\end{lstlisting}

The naming convention is contractual and worth stating explicitly, because
almost every differentiability bug in \cref{sec:ad-rules} is a place where the
wrong member of a pair was chosen:

\begin{table}[H]
  \centering
  \caption{Differentiable and non-differentiable container pairs declared in
    \code{cpp/lupnt/core/definitions.h}.}
  \label{tab:ad-types}
  \begin{tabularx}{\textwidth}{llL}
    \toprule
    Differentiable & Value only & Role \\
    \midrule
    \cpp{Real} & \cpp{double} & Scalar. \cpp{Real} carries one derivative
      slot; \cpp{double} carries none. \\
    \cpp{Vec3}, \cpp{Vec6}, \cpp{VecX} & \cpp{Vec3d}, \cpp{Vec6d},
      \cpp{VecXd} & States, positions, velocities, measurement vectors. \\
    \cpp{Mat3}, \cpp{MatX} & \cpp{Mat3d}, \cpp{MatXd} & Rotations, and the
      output Jacobians. Jacobians are always \cpp{MatXd}: a derivative is a
      number, not a differentiable quantity. \\
    \bottomrule
  \end{tabularx}
\end{table}

\begin{lupntnote}
  The value of a \cpp{Real} is recovered with \code{.val()}; an Eigen
  expression over \cpp{Real} is converted with \code{cast<double>()}. Both
  \emph{discard} the derivative slot. That is exactly what one wants when
  logging, comparing, or handing a result to a non-differentiated interface,
  and exactly what one must not do in the middle of a model
  (\cref{sec:ad-rules}).
\end{lupntnote}

\section{Forward Mode and Why It Was Chosen}
\label{sec:ad-forward}

Let $F : \mathbb{R}^{n} \to \mathbb{R}^{m}$ be the function realised by a
piece of \lupnt{} model code. Forward accumulation evaluates $F$ and,
simultaneously, its directional derivative along a seed vector
$\vec{v} \in \mathbb{R}^{n}$:
\begin{equation}
  \label{eq:ad-jvp}
  \big(F(x),\ J(x)\,\vec{v}\big),
  \qquad J(x) = \frac{\partial F}{\partial x} \in \mathbb{R}^{m\times n},
\end{equation}
at a cost of a small constant multiple of one evaluation of $F$
\citep[Ch.~3]{griewank2008}. Because \cpp{Real} stores a single derivative
slot, one sweep delivers one Jacobian--vector product; a full Jacobian
therefore costs $n$ sweeps, one per input direction:
\begin{equation}
  \label{eq:ad-cost}
  \mathrm{cost}(J) \approx n \cdot c_{F},
  \qquad c_F = \text{cost of one \cpp{Real} evaluation of } F .
\end{equation}
Reverse mode would deliver the same Jacobian in $m$ sweeps, which is cheaper
whenever $m \ll n$. \lupnt{} nonetheless uses forward mode throughout, for
three reasons that are properties of the problem rather than of the library:

\begin{enumerate}[leftmargin=1.5cm]
  \item The dominant object is the state-transition matrix, for which
    $m = n$ (typically $6$ to $12$), so neither mode has an asymptotic
    advantage.
  \item Reverse mode requires taping the entire execution, which for a
    multi-day propagation with a $70\times70$ gravity field and SPICE
    ephemeris lookups is a prohibitive memory cost.
  \item Forward mode is non-invasive: it requires only that the code be
    templated on, or written in, the scalar type. \lupnt{}'s force models,
    frame rotations, time conversions, and measurement models are all
    ordinary C++ written in \cpp{Real}, with no differentiation-specific
    structure.
\end{enumerate}

The order is likewise first only. \code{autodiff::real} is
\code{Real<1,double>}; a Hessian would require \code{Real<2,double>} and is
not instantiated anywhere in \lupnt{}. Every $\Phi$ and every $H$ in this
manual is therefore a first-order object, and second-order filters are out of
scope.

\section{Seeding and Derivative Extraction}
\label{sec:ad-api}

\subsection{The three primitives}

\lupnt{} imports exactly four names from \code{autodiff} into its own
namespace (\cref{lst:ad-defs}). Three of them form the API:

\begin{description}[leftmargin=1.8cm,style=nextline]
  \item[\code{wrt(...)}] Names the independent variables. It may name a whole
    vector (\code{wrt(x)}), a single component (\code{wrt(x(i))}), or a
    separate parameter block (\code{wrt(p)}). These are the variables that
    receive a unit seed in slot $1$.
  \item[\code{at(...)}] Supplies the full evaluation point, in the order the
    function's parameters are declared. It may contain arguments that are
    \emph{not} being differentiated --- times, controls, noise covariances,
    pointers --- and those are simply passed through.
  \item[\code{jacobian(f, wrt(...), at(...), y, J)}] Evaluates $y = f(\cdot)$
    and fills $J$ with $\partial y / \partial(\text{wrt variables})$ in one
    call. It loops internally over the seeded variables, performing one
    forward sweep per variable and writing the resulting column of $J$.
\end{description}

The fourth, \code{derivative(f, wrt(t), at(t))}, returns a single directional
derivative rather than a matrix and is used where the independent variable is
scalar. \cpp{FixedPointingDynamics} uses it to obtain the time derivative of a
rotation matrix, from which the angular velocity follows:

\begin{lstlisting}[language=cpp, caption={Scalar directional derivative: $\dot{R}$ from $R(t)$ by AD in time. Source: \code{cpp/lupnt/dynamics/attitude\_dynamics.cc}}, label={lst:ad-derivative}]
// cpp/lupnt/dynamics/attitude_dynamics.cc :: FixedPointingDynamics
Real t = tf.val();
Mat3 R_b2w = func(tf).reshaped(3, 3);
Mat3 R_b2w_dot = derivative(func, wrt(t), at(t)).reshaped(3, 3);
Vec3 w_w = RotToAngularVelocity(R_b2w, R_b2w_dot);
Vec4 q_b2w = RotToQuat(R_b2w);
\end{lstlisting}

\subsection{The seed-reset idiom}
\label{sec:ad-seedreset}

A \cpp{Real} that arrives from an enclosing computation may already carry a
non-zero derivative slot --- it was itself seeded, or it is downstream of
something that was. Differentiating such a value again would compose the two
seeds and produce the derivative with respect to the \emph{outer} variable,
not the one requested. Every independent \code{jacobian} call in \lupnt{}
therefore begins by copying its input through \code{cast<double>()}, which
constructs a fresh \cpp{VecX} whose derivative slots are all zero:
\begin{equation}
  \label{eq:ad-seedreset}
  \tilde{x} = (x, \dot{x})
  \ \xrightarrow{\ \code{cast<double>()}\ }\ x
  \ \xrightarrow{\ \text{assign to } \cpp{VecX}\ }\ (x, 0),
\end{equation}
after which \code{wrt} plants a clean unit seed. The idiom is

\begin{lstlisting}[language=cpp, caption={Canonical seed reset. This two-line pattern precedes every independent Jacobian in the code base.}, label={lst:ad-idiom}]
VecX x_tmp = x.cast<double>();               // drop any incoming seeds
jacobian(f, wrt(x_tmp), at(x_tmp), y, J);    // J = df/dx evaluated at x
\end{lstlisting}

\begin{lupntwarning}
  Omitting the reset does not raise an error. It produces a Jacobian that is
  numerically plausible but wrong by whatever the stale seed encoded --- most
  often a factor equal to some unrelated sensitivity, or a matrix that is
  correct on the first call and wrong on every subsequent one. Because the
  \code{at(...)} argument must be the same object that \code{wrt(...)} names,
  the reset also fixes the evaluation point: pass \code{x\_tmp}, never the
  original \code{x}, to \code{at}.
\end{lupntwarning}

There is one deliberate exception. \cpp{IslCrosslinkMeasurement::Compute}
re-seeds a fresh state precisely so that the derivative seeds carried by the
\emph{incoming} state are preserved for the consider-state update of the
Schmidt EKF (\cref{chap:filters}); it differentiates a copy and leaves the
caller's seeds untouched.

\begin{lstlisting}[language=cpp, caption={Deliberate fresh re-seed that preserves the caller's derivative directions. Source: \code{cpp/lupnt/measurements/crosslink\_measurement.cc}}, label={lst:ad-isl}]
// cpp/lupnt/measurements/crosslink_measurement.cc :: IslCrosslinkMeasurement::Compute
if (H != nullptr) {
  // Reseed a fresh autodiff state (drops incoming derivative seeds) and differentiate
  // the raw-VecX model, matching the original inline behavior.
  VecX x_tmp = x.cast<double>();
  auto f = [this](const VecX& xx) -> VecX { return Model(xx); };
  jacobian(f, wrt(x_tmp), at(x_tmp), y, *H);
} else {
  y = Model(x);
}
\end{lstlisting}

\section{Jacobian Assembly}
\label{sec:ad-jacobian}

\lupnt{} assembles Jacobians in two ways, differing only in how the $n$
forward sweeps of \cref{eq:ad-cost} are organised.

\subsection{Batched: one \code{jacobian} call}

The default is to hand the whole input vector to \code{jacobian} and let the
library iterate:
\begin{equation}
  \label{eq:ad-batched}
  \big(y,\ J\big) = \code{jacobian}\!\left(f,\ \code{wrt}(x),\ \code{at}(x)\right),
  \qquad
  J_{ji} = \frac{\partial y_j}{\partial x_i}.
\end{equation}
This is used wherever the sweeps are cheap relative to the surrounding work
--- measurement models, analytic dynamics, frame conversions.

\subsection{Column-parallel: \cpp{JacobianParallel}}
\label{sec:ad-jacparallel}

When one sweep means re-propagating an entire trajectory
(\cref{sec:int-stm}), the $n$ sweeps dominate. Since the columns of $J$ are
mutually independent,
\begin{equation}
  \label{eq:ad-col}
  J\,\vec{e}_i = \left.\frac{\partial}{\partial\epsilon}\right|_{\epsilon=0}
    F\!\left(x + \epsilon\,\vec{e}_i\right),
  \qquad i = 1,\dots,n,
\end{equation}
\cpp{JacobianParallel} spawns one OpenMP thread per column, each seeding a
single component \code{wrt(x0\_tmp(i))} and running its own sweep, and then
assembles the $m\times n$ result. Note that each thread performs its own seed
reset, which is what makes the columns independent and the loop thread-safe
with respect to the AD state.

\begin{lstlisting}[language=cpp, caption={Column-parallel Jacobian. Source: \code{cpp/lupnt/numerics/math\_utils.cc}}, label={lst:ad-jacpar}]
// cpp/lupnt/numerics/math_utils.cc :: JacobianParallel
VecX JacobianParallel(const std::function<VecX(const VecX&)>& func, const VecX& x0, MatXd& J) {
  int n = x0.size();
  std::vector<MatXd> J_vec(n);
  VecX xf;

#pragma omp parallel for
  for (int i = 0; i < n; i++) {
    VecX xi;
    MatXd Ji;
    VecX x0_tmp = x0.cast<double>();
    jacobian(func, wrt(x0_tmp(i)), at(x0_tmp), xi, Ji);
    if (i == 0) xf = xi;
    J_vec[i] = Ji;
  }

  int m = xf.size();
  J.resize(m, n);
  for (int i = 0; i < n; i++) {
    for (int j = 0; j < m; j++) {
      J(j, i) = J_vec[i](j);
    }
  }
  return xf;
}
\end{lstlisting}

\begin{lupntnote}
  \cpp{JacobianParallel} returns the function value from the $i=0$ thread
  only. The value is the same in every thread --- forward-mode AD does not
  perturb the value slot --- so this is a valid, if fragile, optimisation. It
  does mean that a function with $n = 0$ inputs leaves \code{xf} empty and
  \code{J} unsized.
\end{lupntnote}

\section{State-Transition and Parameter Sensitivity Matrices}
\label{sec:ad-stm}

\subsection{State-transition matrix}

The default \cpp{Dynamics::Propagate} with a non-null \code{stm} pointer
differentiates the seed-reset, no-STM propagation with respect to the initial
state:
\begin{equation}
  \label{eq:ad-stm}
  \Phi(t_f, t_0) = \frac{\partial x(t_f)}{\partial x(t_0)}
    = \code{jacobian}\!\Big(
        x_0 \mapsto \code{Propagate}(x_0, t_0, t_f, u),
        \ \code{wrt}(x_0),\ \code{at}(x_0)
      \Big).
\end{equation}
This is the same mathematical object as the integrator sensitivity of
\cref{eq:int-jac}; the two differ only in batching (one \code{jacobian} call
versus \cpp{JacobianParallel}). In both cases what is differentiated is the
\emph{discrete} propagation map, so $\Phi$ inherits the integrator's
truncation error and is consistent with the trajectory that is actually
produced --- the property a covariance propagation requires. The continuous
variational contract $\dot{\Phi} = A(t)\Phi$ with
$A = \partial f/\partial x$ is documented in \cref{chap:dynamics}; it is never
integrated.

\subsection{Joint state and parameter sensitivities}

Batch and EKF estimation of force-model or clock parameters needs the two
blocks
\begin{equation}
  \label{eq:ad-stmparam}
  \Phi_{x} = \frac{\partial x(t_f)}{\partial x(t_0)},
  \qquad
  \Phi_{p} = \frac{\partial x(t_f)}{\partial p},
\end{equation}
obtained by seeding the two argument blocks in turn while holding the
evaluation point fixed. The function is evaluated twice, once per seeded
block; \code{at(x0\_tmp, p0\_tmp)} is identical in both calls, so both
derivatives are taken at the same trajectory.

\begin{lstlisting}[language=cpp, caption={STM and parameter sensitivity by seeded blocks. Source: \code{cpp/lupnt/dynamics/dynamics.cc}}, label={lst:ad-stm}]
// cpp/lupnt/dynamics/dynamics.cc :: Dynamics::Propagate(..., MatXd* stm)
auto func = [=, this](const VecX& x) {
  State xs = State(x, x0.GetNames(), x0.GetUnits());
  return Propagate(xs, t0, tf, u);
};
VecX xf_tmp;
VecX x0_tmp = x0.cast<double>();
jacobian(func, wrt(x0_tmp), at(x0_tmp), xf_tmp, *stm);       // stm = d(xf)/d(x0)

// cpp/lupnt/dynamics/dynamics.cc :: Dynamics::PropagateWithParams(..., stm_state, stm_param)
VecX x0_tmp = x0.cast<double>();
VecX p0_tmp = params.cast<double>();
jacobian(func, wrt(x0_tmp), at(x0_tmp, p0_tmp), xf_tmp, *stm_state);  // d(xf)/d(x0)
jacobian(func, wrt(p0_tmp), at(x0_tmp, p0_tmp), xf_tmp, *stm_param);  // d(xf)/d(p)
\end{lstlisting}

\cpp{NBodyDynamics} overrides the STM overload with a hard precondition:
differentiation is refused unless AD has been enabled, because its gravity
field has a \cpp{double} fast path that is not differentiable
(\cref{sec:ad-rules}).

\begin{lstlisting}[language=cpp, caption={AD precondition on the N-body STM. Source: \code{cpp/lupnt/dynamics/numerical\_orbit\_dynamics.cc}}, label={lst:ad-nbody}]
// cpp/lupnt/dynamics/numerical_orbit_dynamics.cc :: NBodyDynamics::Propagate
State NBodyDynamics::Propagate(const State& x0, Real t0, Real tf, const State* u, MatXd* stm) {
  LUPNT_CHECK(use_ad_, "Autodiff not enabled", "NBodyDynamics");
  (void)u;
  return NumericalDynamics::Propagate(x0, t0, tf, u, stm);
}
\end{lstlisting}

\section{Coupling to the Filters}
\label{sec:ad-filters}

The estimators of \cref{chap:filters} require, at each step, a linearisation
alongside the nonlinear evaluation:
\begin{equation}
  \label{eq:ad-fh}
  F = \frac{\partial x(t_f)}{\partial x(t_0)} \in \mathbb{R}^{n\times n},
  \qquad
  H = \frac{\partial z}{\partial x} \in \mathbb{R}^{n_z \times n},
\end{equation}
$F$ for the covariance prediction $P^- = F P F\transpose + Q$ and $H$ for the
innovation covariance and gain. \lupnt{} does not ask model authors to supply
either. A plain \cpp{DynamicsFunction} $(x, t_0, t_f, u) \mapsto x(t_f)$ and a
plain \cpp{MeasurementFunction} $(x, R) \mapsto z$ are \emph{lifted} into the
filter signatures \cpp{FilterDynamicsFunction} and
\cpp{FilterMeasurementFunction} by wrapping them in \code{jacobian}. This is
the single point at which AD enters the estimation stack.

\begin{lstlisting}[language=cpp, caption={Lifting plain models into filter callbacks. Source: \code{cpp/lupnt/numerics/filters/filter\_utils.cc}}, label={lst:ad-filter}]
// cpp/lupnt/numerics/filters/filter_utils.cc :: GetFilterMeasurementFunction
FilterMeasurementFunction f_meas_func = [f_meas](const State& x, MatXd* H, MatXd* R) -> VecX {
  VecX y;
  VecX x_tmp = x.cast<double>();
  jacobian(f_meas, wrt(x_tmp), at(x_tmp, R), y, *H);      // H = dz/dx
  return y;
};

// cpp/lupnt/numerics/filters/filter_utils.cc :: GetFilterDynamicsFunction
FilterDynamicsFunction f_dyn_func
    = [f_dyn](const State& x, Real t0, Real tf, const State* u, MatXd* F) -> State {
  VecX xf;
  VecX x_tmp = x.cast<double>();
  jacobian(f_dyn, wrt(x_tmp), at(x_tmp, t0, tf, u), xf, *F);   // F = d(xf)/d(x0)
  return State(xf, x.GetNames(), x.GetUnits());
};
\end{lstlisting}

Three consequences follow, and they are the reason \cref{sec:ad-rules} is a
requirement rather than a style guide.

\begin{enumerate}[leftmargin=1.5cm]
  \item \textbf{The filter's linearisation is the model's derivative.} There
    is no separate, independently maintained Jacobian that could disagree with
    the model. A change to a force model automatically changes $F$.
  \item \textbf{A non-differentiable model yields a structurally wrong $F$ or
    $H$, not an exception.} A block of the state that a model handles in
    \cpp{double} arithmetic contributes an exactly zero block. The filter then
    believes that part of the state is unobservable or unaffected by dynamics,
    and the covariance becomes optimistic in a way no residual test detects.
  \item \textbf{The UKF is the exception.} \cpp{UKF::Predict} calls the
    dynamics callback with \code{F = nullptr}, because it propagates sigma
    points and needs no linearisation \citep{julier2004ukf}. A model that is
    non-differentiable but numerically correct will therefore work in the UKF
    and fail silently in the EKF, SRIF, and UDU filters
    \citep{bierman1977, thornton1980}.
\end{enumerate}

\section{Differentiability Rules for Models}
\label{sec:ad-rules}

A model is differentiable in the sense required here if, along the execution
path actually taken, every operation from input to output is one of the
overloaded \cpp{Real} operations of
\cref{eq:ad-arith}--\cref{eq:ad-chain} with a finite derivative factor. The
following rules are what that means in practice.

\subsection{Rule 1: never narrow to \cpp{double} inside a model}

Assigning a \cpp{Real} to a \cpp{double}, or an Eigen \cpp{Real} expression to
a \code{d}-suffixed container, drops slot $1$. If the value is later widened
back to \cpp{Real} it re-enters with a zero derivative, and the whole
downstream Jacobian block is zero. The canonical example is the gravity-field
fast path in \cpp{NBodyDynamics::ComputeRates}, which is exactly why the STM
overload of \cref{lst:ad-nbody} refuses to run with \code{use\_ad\_} disabled.

\begin{lstlisting}[language=cpp, caption={Rule 1. The \code{else} branch is numerically correct and produces a zero derivative block. Source: \code{cpp/lupnt/dynamics/numerical\_orbit\_dynamics.cc}}, label={lst:ad-rule1}]
// cpp/lupnt/dynamics/numerical_orbit_dynamics.cc :: NBodyDynamics::ComputeRates
if (use_ad_) {
  // DO: stay in Real -- the spherical-harmonic sum is differentiated exactly
  a_bf = AccelarationGravityField<Real>(r_bf, grav.GM, grav.R, grav.CS, grav.n, grav.m);
} else {
  // DON'T (for any STM/filter use): every derivative is discarded here
  Vec3d r_bf_d = r_bf.cast<double>();
  a_bf = AccelarationGravityField<double>(r_bf_d, static_cast<double>(grav.GM),
                                          static_cast<double>(grav.R),
                                          grav.CS.cast<double>(), grav.n, grav.m);
}
\end{lstlisting}

Use \code{.val()} and \code{cast<double>()} only at the boundaries: logging,
formatted output, comparisons that select a branch, storage into a results
container, and the deliberate seed reset of \cref{sec:ad-seedreset}.

\subsection{Rule 2: branch on values, but understand what the branch means}

Comparing \cpp{Real}s is unavoidable and harmless in itself --- comparison
operators read slot $0$ only. What matters is that \code{jacobian} returns the
derivative of the branch \emph{actually taken}. Two acceptable patterns:

\begin{lstlisting}[language=cpp, caption={Rule 2, acceptable. Selection returns one of the operands unchanged, so its derivative travels with it. Source: \code{cpp/lupnt/numerics/math\_utils.cc}}, label={lst:ad-rule2ok}]
// cpp/lupnt/numerics/math_utils.cc
Real Max(Real x, Real y) { return x.val() > y.val() ? x : y; }
Real Min(Real x, Real y) { return x.val() < y.val() ? x : y; }

// cpp/lupnt/numerics/math_utils.cc :: safe_acos -- round-off guard, not a model branch
Real safe_acos(Real x) {
  if (x >= 1.0) {
    return acos(x - EPS);
  } else if (x <= -1.0) {
    return acos(x + EPS);
  } else {
    return acos(x);
  }
}
\end{lstlisting}

\cpp{Max} and \cpp{Min} select a whole \cpp{Real} and therefore carry its
derivative; they compute a valid one-sided derivative at the kink. What is
\emph{not} acceptable is a branch that reconstructs a value from
\code{.val()}, because that rebuilds the number with a zero seed:

\begin{lstlisting}[language=cpp, caption={Rule 2, the failure mode. Both branches return a derivative-free constant.}, label={lst:ad-rule2bad}]
// DON'T
Real ClampBad(Real x, double lo, double hi) {
  if (x.val() < lo) return Real(lo);          // derivative 0 -- correct here
  if (x.val() > hi) return Real(hi);          // derivative 0 -- correct here
  return Real(x.val());                       // derivative 0 -- WRONG: x was interior
}

// DO
Real ClampGood(Real x, double lo, double hi) {
  if (x.val() < lo) return Real(lo);
  if (x.val() > hi) return Real(hi);
  return x;                                   // interior: pass the Real through
}
\end{lstlisting}

\subsection{Rule 3: keep away from the singular edges of the elementary
  functions}

By \cref{eq:ad-acos}, $\arcsin$ and $\arccos$ have infinite derivative at
$|x| = 1$ and produce $\mathrm{NaN}$ there. The solar-radiation-pressure
shadow function evaluates all three of $\arcsin(R_\odot/\rho)$,
$\arcsin(R_b/r)$, and $\arccos(\cdot)$ on quantities that saturate exactly at
a grazing eclipse boundary. \cpp{ClampUnit} returns a value strictly inside
the domain, $\pm(1 - 10^{-12})$, rather than the endpoint: the value shift is
below \num{1e-9} and irrelevant, while the derivative becomes finite (in fact
zero, since the returned constant no longer depends on $x$). This is what
allows a state-transition matrix to be formed across an arc that grazes a
shadow boundary.

\begin{lstlisting}[language=cpp, caption={Rule 3. Clamp to the interior, never to the endpoint. Source: \code{cpp/lupnt/environment/forces.cc}}, label={lst:ad-clampunit}]
// cpp/lupnt/environment/forces.cc :: ClampUnit
Real ClampUnit(Real x) {
  // Clamp to [-1, 1] for the asin/acos calls in ShadowFunction. When the argument
  // saturates we return a constant *strictly inside* the domain (+/-(1 - kEps))
  // rather than exactly +/-1: under autodiff asin'/acos' diverge at +/-1, and
  // returning exactly +/-1 yields 0/0 = NaN in the derivative (the value is
  // unaffected). Returning a strictly-interior constant keeps the derivative
  // finite (it is zero there, since the return no longer depends on x). This is
  // what lets the state-transition matrix be formed analytically over arcs that
  // graze a shadow boundary. The value shift is < 1e-9 and irrelevant.
  constexpr double kEps = 1e-12;
  if (x.val() >= 1.0) return 1.0 - kEps;
  if (x.val() <= -1.0) return -1.0 + kEps;
  return x;
}
\end{lstlisting}

The same reasoning applies to $\sqrt{\tilde{x}}$ at $x = 0$ (see
\cref{eq:ad-arith}), which is why \file{cpp/lupnt/environment/forces.cc}
writes \code{sqrt(Max(y2, 0.0))} rather than \code{sqrt(y2)}, and to
$\|\vec{v}\|$ at $\vec{v} = \vec{0}$.

\subsection{Rule 4: build special functions from arithmetic and \code{sqrt}}
\label{sec:ad-agm}

The C++ standard library's special functions --- \code{std::comp\_ellint\_1},
\code{std::cyl\_bessel\_j}, and the rest --- take and return \cpp{double}.
Calling one inside a model violates Rule 1 twice over: it narrows the input
and returns a seedless output. The remedy is to compose the function out of
operations that \cpp{Real} already overloads.

The complete elliptic integral of the first kind, needed for the ring
potential that represents the DE440 asteroid and Kuiper populations in the
relativistic time transformation of \cref{chap:time-conversions}, is the
worked example. The arithmetic--geometric mean of $a_0$ and $b_0$ is the
common limit of
\begin{equation}
  \label{eq:ad-agm}
  a_{n+1} = \tfrac{1}{2}\left(a_n + b_n\right),
  \qquad
  b_{n+1} = \sqrt{a_n b_n},
\end{equation}
which converges quadratically --- the number of correct digits doubles each
iteration, so roughly five iterations reach double precision --- and gives
\citep[\S17.6]{abramowitz1964}
\begin{equation}
  \label{eq:ad-agm-K}
  K(k) = \frac{\pi}{2\,\mathrm{AGM}\!\left(1,\ k'\right)},
  \qquad k' = \sqrt{1-k^{2}} .
\end{equation}
Every operation in \cref{eq:ad-agm}--\cref{eq:ad-agm-K} is an addition, a
multiplication, a division, or a square root, all of which appear in
\cref{eq:ad-arith}. The iteration is therefore differentiable as written,
without any special-function derivative rule, and the derivative converges at
the same quadratic rate as the value.

\begin{lstlisting}[language=cpp, caption={Rule 4. An AGM iteration replaces a non-differentiable library special function. Source: \code{cpp/lupnt/conversions/time\_conversions.cc}}, label={lst:ad-agm}]
// cpp/lupnt/conversions/time_conversions.cc :: RingPotential
/// K is evaluated by the arithmetic-geometric mean, K(k) = pi/(2 AGM(1,k')),
/// k' = sqrt(1-k^2). AGM converges quadratically (~5 iterations here) and is
/// differentiable, so this composes with autodiff unlike std::comp_ellint_1.
Real RingPotential(Real r, double gm_ring, double a_ring) {
  if (gm_ring == 0.0) return Real(0.0);
  Real apr = a_ring + r;
  Real k2 = 4.0 * a_ring * r / (apr * apr);
  Real a = 1.0;
  Real b = sqrt(1.0 - k2);          // complementary modulus k'
  for (int i = 0; i < 12; i++) {
    Real a_next = 0.5 * (a + b);
    b = sqrt(a * b);
    a = a_next;
    if (abs((a - b).val()) < 1e-16) break;
  }
  Real K = PI / (2.0 * a);
  return 2.0 * gm_ring * K / (PI * apr);
}
\end{lstlisting}

Note the convergence test: it compares \code{.val()} only, so the loop trip
count is decided by values and the derivative simply rides along the fixed
number of iterations actually executed. That is the correct pattern for an
iterative algorithm --- differentiate the iterates, do not differentiate the
stopping rule.

\subsection{Rule 5: know which helpers deliberately pass the seed through}

\lupnt{}'s rounding helpers overwrite \emph{only} slot $0$ and leave slot $1$
untouched:

\begin{lstlisting}[language=cpp, caption={Rule 5. Value-modifying, seed-preserving helpers. Source: \code{cpp/lupnt/numerics/math\_utils.cc}}, label={lst:ad-rule5}]
// cpp/lupnt/numerics/math_utils.cc
Real floor(Real x) { Real y = x; y[0] = std::floor(x.val()); return y; }
Real ceil(Real x)  { Real y = x; y[0] = std::ceil(x.val());  return y; }
Real frac(Real x)  { Real y = x; y[0] = x.val() - std::floor(x.val()); return y; }
Real mod(Real x, Real y) { Real z = x; z[0] = std::fmod(x.val(), y.val()); return z; }
Real round(Real x, int n) {
  Real y = x;
  y[0] = std::round(x.val() * std::pow(10, n)) / std::pow(10, n);
  return y;
}
\end{lstlisting}

Mathematically $\lfloor x \rfloor$ has derivative $0$ almost everywhere, so
these helpers do \emph{not} return the true derivative of the function they
name. The behaviour is intentional and is what the time-scale conversions of
\cref{chap:time-conversions} require: splitting an epoch as
$t = \lfloor t \rfloor + \mathrm{frac}(t)$ must preserve
$\partial t/\partial\,(\cdot) = 1$ across the split, which it does precisely
because \code{floor} and \code{frac} each pass the seed through and the
derivative of the sum is recovered. Used anywhere else --- a genuine
quantisation, a table index, a modulo wrap whose discontinuity is physical ---
they will report a derivative of $1$ where the true value is $0$. Check the
intent before reaching for them.

\subsection{Checklist}

\begin{table}[H]
  \centering
  \caption{Differentiability checklist for a new model.}
  \label{tab:ad-checklist}
  \begin{tabularx}{\textwidth}{LL}
    \toprule
    Do & Don't \\
    \midrule
    Write the model in \cpp{Real}; template on the scalar if a
      \cpp{double} fast path is also needed. &
    Assign intermediates to \cpp{double} / \code{Vec3d} / \code{MatXd}
      inside the model. \\
    Reset seeds with \code{cast<double>()} before an independent
      \code{jacobian}. &
    Pass the original \code{x} to \code{at(...)} while naming a copy in
      \code{wrt(...)}. \\
    Clamp \code{asin}/\code{acos} arguments to the strict interior
      (\cpp{ClampUnit}). &
    Clamp them to exactly $\pm 1$, or call \code{acos} on an unguarded
      dot product. \\
    Guard \code{sqrt} and \code{norm} against a zero argument with
      \cpp{Max}. &
    Take \code{sqrt} of a quantity that can round to a small negative
      number. \\
    Compose special functions from $+$, $-$, $\times$, $\div$,
      \code{sqrt} (AGM, series). &
    Call \code{std::comp\_ellint\_1} or any \cpp{double}-only special
      function. \\
    Branch on \code{.val()} and return the selected \cpp{Real}
      unchanged. &
    Rebuild the return value as \code{Real(x.val())}. \\
    Use \code{.val()} at output boundaries only. &
    Use \cpp{floor}/\cpp{round}/\cpp{mod} where a true zero derivative
      is meant. \\
    \bottomrule
  \end{tabularx}
\end{table}

\section{Model Boundaries}
\label{sec:ad-boundaries}

\paragraph{First order, forward mode only.}
\cpp{Real} is \code{Real<1,double>}. \lupnt{} produces Jacobians and nothing
else: no Hessians, no second-order filters, no curvature-based step
selection. A full $m\times n$ Jacobian costs $n$ forward sweeps
(\cref{eq:ad-cost}); there is no reverse-mode path, so a scalar objective over
a large parameter vector --- a cost function for trajectory optimisation, for
instance --- is differentiated at a cost proportional to the parameter count
rather than at constant cost.

\paragraph{The derivative is of the code, not of the mathematics.}
AD differentiates the execution path actually taken. Where a model clamps, a
table is indexed, an iteration terminates, or a branch is selected, the
reported derivative is that of the realised path. \cpp{ClampUnit} returns a
derivative of zero at a saturated shadow boundary, which is finite and useful
but is not the derivative of the true, non-differentiable shadow function.
Similarly, an ephemeris read from a tabulated kernel contributes the
derivative of the interpolant, not of the underlying body motion.

\paragraph{Seed hygiene is unenforced.}
The \code{cast<double>()} reset of \cref{eq:ad-seedreset} is a convention, not
a type-system guarantee. Omitting it compiles, runs, and yields a plausible
but incorrect Jacobian. There is no runtime check that an input arrived
unseeded.

\paragraph{Silent zero blocks.}
Any part of a model written in \cpp{double} contributes an exactly zero block
to $F$ or $H$ rather than an error. \cpp{NBodyDynamics} guards its own case
with \code{LUPNT\_CHECK(use\_ad\_, ...)} (\cref{lst:ad-nbody}), but that guard
is specific to that class; nothing prevents another model from silently
returning a zero-derivative block. The recommended acceptance test for a new
model is to compare its AD Jacobian against a central finite difference of its
own value function, component by component, at a few representative states.

\paragraph{Cost and memory.}
A \cpp{Real} is twice the size of a \cpp{double} and roughly doubles the
arithmetic of every operation, so a \cpp{Real} evaluation of a force model
costs about twice its \cpp{double} equivalent even when no derivative is
wanted. Combined with \cref{eq:ad-cost}, a six-state STM over a trajectory
therefore costs on the order of $12$ times a plain \cpp{double} propagation,
reduced by the OpenMP fan-out of \cpp{JacobianParallel}
(\cref{sec:ad-jacparallel}).

\paragraph{Thread safety.}
\cpp{JacobianParallel} is safe because each thread owns its own seeded copy of
the input. That guarantee does not extend to whatever the differentiated
function touches: a model that reads or writes global mutable state --- the
global reference epoch being the notable example --- is not automatically safe
to differentiate in parallel.

\paragraph{Accuracy relative to finite differences.}
AD carries no truncation error and no subtractive cancellation: the result is
the derivative of the floating-point program to within round-off
\citep{ieee754}, with no step-size parameter to tune. Where documentation
elsewhere in the repository describes a Jacobian as ``finite-difference
based'' --- including the Doxygen comment on the integrator's sensitivity
overload (\cref{sec:int-stm}) --- that comment is stale; the implementation
calls \code{autodiff::jacobian}.
